\subsection{Single Pole Active Low Pass}\label{sec:filt_active_sp}

\subsubsection{Overview}
An inverting amplifier (Section~\ref{sec:amp_inverting}) with \(C\) across
the feedback resistor: an RC low pass with gain and a buffered output.

\subsubsection{Schematic}
\begin{circuitfig}
  \draw (0,0) node[op amp](A){};
  \draw (A.-) -- ++(-0.6,0) coordinate(n) node[circ]{} node[below left]{$V_-$} to[european resistor, l_={$G_1$}] ++(-2.4,0) node[circ]{} node[left]{$V_{in}$};
  \draw (n) |- (-0.5,1.6) to[european resistor, l={$G_2$}] (1.6,1.6) -| (A.out);
  \draw (-0.5,1.6) |- (-0.5,2.8) to[capacitor, l={$C$}] (1.6,2.8) -- (1.6,1.6);
  \draw (A.+) -- ++(-0.6,0) node[ground]{};
  \draw (A.out) to[short, -o] ++(0.8,0) node[right]{$V_{out}$};
\end{circuitfig}

\subsubsection{Transfer Function}
\textbf{Small Signal, Frequency Domain, AC Input Analysis (Q4)} with an ideal op-amp (at DC
\(C\) is open and this is the inverting amplifier, Section~\ref{sec:amp_inverting}): \(V_+ = 0\),
so \(V_- = 0\) and no current enters \(V_-\); the feedback admittance is \(G_2 + sC\).
Applying KCL at \(V_-\):
\begin{align*}
  (V_- - V_{in}) G_1 + (V_- - V_{out})(G_2 + sC) &= 0 \\
  -V_{in} G_1 - V_{out}(G_2 + sC) &= 0 \hspace{1cm} [V_- = 0] \\
  H(s) = \frac{V_{out}}{V_{in}} &= -\frac{G_1}{G_2 + sC} = -\frac{\frac{1}{R_1}}{\frac{1 + sR_2 C}{R_2}} \hspace{1cm} \left[G_k = \tfrac{1}{R_k}\right] \\
  H(s) &= -\frac{R_2}{R_1}\cdot\frac{1}{1 + sR_2 C} \tag{1}
\end{align*}

\subsubsection{Analysis}
\begin{enumerate}
  \item \textbf{Poles (\(D(s) = 0\))}:
    \begin{align*}
      1 + s_p R_2 C &= 0 \hspace{1cm} \Rightarrow \hspace{1cm} s_p = -\frac{1}{R_2 C} \hspace{1cm} f_p = \frac{|s_p|}{2\pi} = \frac{1}{2\pi R_2 C}
    \end{align*}
  \item \textbf{Zeroes (\(N(s) = 0\))}: \(N(s) = -\frac{R_2}{R_1} \neq 0\): no finite zero.
  \item \textbf{DC Gain (\(H(0)\))}: \(s = 0\) in (1), \(H(0) = -\frac{R_2}{R_1}\cdot\frac{1}{1 + 0} = -\frac{R_2}{R_1}\).
  \item \textbf{Gain-Bandwidth Product (\(GBW\))}:
    \(GBW = |H(0)| f_p = \frac{R_2}{R_1}\cdot\frac{1}{2\pi R_2 C} = \frac{1}{2\pi R_1 C}\), set by \(R_1\) and \(C\)
    only (the op-amp needs \(GB \gg GBW\)).
  \item \textbf{Impedances}:
    \paragraph{At Input:} \(V_- = 0\) (virtual ground):
    \(Z_{in} = \frac{V_{in}}{(V_{in} - 0) G_1} = R_1\).
    \paragraph{At Output:} the op-amp output, \(Z_{out} = \frac{R_o}{1 + AF} \approx 0\)
    (Section~\ref{sec:amp_inverting}): unlike the passive RC (Section~\ref{sec:theory_general_loading})
    the next stage does not load it.
  \item \textbf{Advantages \& Limitations}: gain and pole set independently
    (\(\frac{R_2}{R_1}\), \(\frac{1}{2\pi R_2 C}\)), buffered output; but inverting, \(Z_{in} = R_1\), only
    \(-20\,\)dB/decade, and an op-amp (power, \(GB\), swing).
\end{enumerate}

\subsubsection{Design}
\begin{enumerate}
  \item \(R_1 = Z_{in}\) required; \(R_2 = |H(0)| R_1\).
  \item \(C = \frac{1}{2\pi R_2 f_p}\).
\end{enumerate}

\subsubsection{Question}
\(R_1 = 1\,k\Omega\), \(R_2 = 10\,k\Omega\), \(C = 15.9\,nF\). Find the DC gain and
\(f_p\).
\paragraph{Solution:}
\begin{align*}
  H(0) &= -\frac{R_2}{R_1} = -\frac{10\,k\Omega}{1\,k\Omega} = -10 \\
  f_p &= \frac{1}{2\pi R_2 C} = \frac{1}{2\pi \cdot 10^4\,\Omega \cdot 15.9 \times 10^{-9}\,F}
       = \frac{1}{9.990 \times 10^{-4}\,s} = 1.001\,kHz
\end{align*}
